\originalchapter{链构象与弹性}
\originalsection{随机链和熵弹簧}
一条柔性链由 $N$ 段长度 $a$ 的独立、各向同性键组成。端到端向量 $\mathbf R=\sum_i\mathbf a_i$，交叉项平均为零，所以 $\avg{R^2}=Na^2$。当 $N\gg1$、延伸远小于轮廓长 $Na$ 时，各分量的中心极限定理给出高斯构象密度
\[
P_0(\mathbf R)=\left(\frac3{2\pi Na^2}\right)^{3/2}\exp\!\left(-\frac{3R^2}{2Na^2}\right).
\]
取 $-\kb T\ln P_0$ 得到自由能中的二次项 $3\kb TR^2/(2Na^2)$。这个弹簧来自可用构象数减少，刚度随温度增加；它不同于一根金属弹簧的键能伸缩。

拉力 $F\hat z$ 引入能量 $-FR_z$。配平方可得 $\avg{R_z}=Na^2F/(3\kb T)$。若需要接近完全拉直，应从单键角分布计算，得到自由连接链的 Langevin 函数
\[
\frac{\avg{R_z}}{Na}=\coth(\beta Fa)-\frac1{\beta Fa}.
\]
这是对每个键积分 $\int\e^{\beta Fa\cos\theta}\dd\Omega$ 后对 $F$ 求导的结果。小力展开恢复高斯弹簧，大力趋于 $1$，不会像高斯模型那样无限延伸。

\originalsection{排除体积和尺度}
好溶剂中的单体不能互相穿透。用第二维里系数 $v>0$ 表示排斥，大小 $R$ 的链平均单体密度为 $N/R^d$。两体排斥自由能估计为 $\kb T vN^2/R^d$，与熵弹性相加：
\[
\frac{F(R)}{\kb T}\sim\frac{R^2}{Na^2}+\frac{vN^2}{R^d}.
\]
对 $R$ 最小化给 $R\sim N^{3/(d+2)}$；三维指数约 $3/5$，是尺度近似而非精确证明。上临界维数 $d=4$ 以上恢复理想链指数 $1/2$。差溶剂中 $v$ 变负，两体项单独会导致塌缩，必须加入正的三体项稳定有限密度，不能继续用好溶剂公式。

处于拉力 $F$ 下的自回避链可按“张力团块”分组。团块大小 $\xi=\kb T/F$，内部尚不感受拉力，$\xi\sim a g^\nu$；链有 $N/g$ 个团块，大尺度排列沿力方向，所以 $X\sim(N/g)\xi\sim Na(Fa/\kb T)^{1/\nu-1}$。适用窗口是零力链大小远大于 $\xi$ 且 $\xi\gg a$；达到键尺度后必须转向饱和模型。

\originalsection{弯曲链与持久长度}
DNA在几十纳米尺度内不能自由折角。用弧长 $s$、单位切向量 $\mathbf t(s)$ 描述轮廓长 $L$ 的不可伸长链，弯曲能为 $H_b=\kappa\int_0^L|\partial_s\mathbf t|^2\dd s/2$。三维WLC模型定义持久长度 $\ell_p=\kappa/(\kb T)$，切向关联为 $\avg{\mathbf t(s)\cdot\mathbf t(0)}=\e^{-s/\ell_p}$。对两次弧长积分有
\[
\avg{R^2}=2\ell_pL-2\ell_p^2(1-\e^{-L/\ell_p}).
\]
$L\ll\ell_p$ 时近似刚棒，$L\gg\ell_p$ 时等效随机链的 Kuhn 长度是 $2\ell_p$。

在强拉力、弱弯曲近似中，两个横向方向的 Fourier 模有能量 $(\kappa q^2+F)|\theta_q|^2/2$。由等分配积分得到 $1-X/L\simeq\tfrac12\sqrt{\kb T/(F\ell_p)}$。此式只给强拉力渐近；常用的整段插值不是从此渐近式严格推出的精确WLC解。原题中的蛋白诱导局部弯折会提供一个额外、可分离的延伸亏损。
